\section{Confining large common divisors}
\label{p13:conf:section}

At small relative rank, the host count already controls
triple probabilities. At macroscopic rank, many large
prime factors must be shared between the windows.
An interpolation argument confines the offsets of these
common divisors to a fixed-degree curve. Unless the pair
has an exact rational alignment, that curve contains too
few offsets to sustain many large groups.

Fix \(H=C_1B\) larger than the diameters and use the odd
kernel \(K_H\) from Section~\ref{p13:half:section}.
For this section, a pair is called aligned if two distinct
offset pairs give the same ratio of true values.
Equivalently, two points \((i,j)\) lie on a line through
\((-x,-y)\). Reducing the ratio bounds its height by
\(O(B)\); the enlarged height-\(B^3\) convention used
elsewhere in the paper includes these alignments.


The confinement lemma, proved in Section~\ref{p14:comp:confinement},
starts from pairwise coprime groups $g_a>N^\tau$ dividing
$x+i_a$ and $y+j_a$.  Determinants of bounded-degree monomials force
their offset points onto an algebraic curve.  Nonlinear components
have $o(B)$ grid points, while a line containing many groups must be
an exact alignment.  After removing these few groups, the surviving
coordinates have small odd kernels and the half-count machinery makes
their centres rare.  The argument never enumerates the auxiliary curve
or a chosen relation subspace.

For each fixed arity, we also use an anchored version:
after coordinate regularization, each active anchor value meets only
a fixed number of minimal prime supports. Confinement removes the
large supports at a cost of $o(B)$ anchor coordinates. Two remaining
small kernels make that centre rare. On a primitive family this
works at every window, giving $N^\varepsilon$ dense hosts and the
fixed-size weights in Proposition~\ref{p13:transfer:primitive}.
The arity and density proportion are fixed before $N$.
The same section completes the stronger pair profile of
Proposition~\ref{p1:arith:hosts} without using that profile as an input.

The two forms needed to partition compatible triples are as follows.
The argument also proves the primitive-triple and dense-pair versions,
Theorem~\ref{p13:conf:dense-primitive} and
Corollary~\ref{p13:conf:dense-pairs}.

\begin{proposition}[Dense triples with weak proper intersections]
\label{p13:conf:near-primitive}
Fix \(\eta>0\) and \(0<\theta<\eta/100\).
Admissible triples satisfying
\[
 \rho_{\mathrm{rel}}(T)\ge\eta B,\qquad
 d_{xy}\le\theta B\quad\text{for every pair}
\]
have \(O_{\eta,\theta,C,\varepsilon}(N^\varepsilon)\)
hosts and total homogeneous mass
\(O(N^{-2+\theta+\varepsilon})\).
\end{proposition}

Its proof is given in Section~\ref{p14:comp:confinement:near}.

\begin{proposition}[A strong nonaligned pair in a triple]
\label{p13:conf:strong-profile}
Fix \(\theta,\sigma,\varepsilon>0\).
Compatible connected triples with a nonaligned pair
of rank at least \(\theta L\), and
\(\rho_{\mathrm{rel}}(T)\le\sigma L\), have total
joint mass
\[
 O_{\theta,\sigma,C,\varepsilon}
       \bigl(N^{-\min\{3/2,\,3-\sigma\}+\varepsilon}\bigr).
\]
\end{proposition}

Its proof is given in Section~\ref{p14:comp:confinement:strong}.



The last proposition preserves the actual conditional
probabilities. Replacing them by independent weights
would lose the large-rank case. What remains is a
pointwise bound on the possible pair ranks and a separate
calculation for rationally aligned triples.

\subsection{Polynomial confinement and dense arithmetic profiles}
\label{p14:comp:confinement}

We use the block and alignment conventions of
Section~\ref{p13:conf:section}. The group argument treats nonaligned
pairs and primitive families at every fixed arity, then proves the
two triple interfaces stated there.

\begin{lemma}[Polynomial confinement of common divisors]
\label{p13:conf:groups}
Fix \(\tau>0\), \(0<\delta<1/2\), and \(C_0>0\).
Let \(x,y\asymp N\), and let pairwise coprime integers
\(g_a>N^\tau\) satisfy
\[
 g_a\mid x+i_a,\quad g_a\mid y+j_a,\qquad
 |i_a|,|j_a|\le C_0B.
\]
Suppose no line through \((-x,-y)\) contains two
distinct points among the \((i_a,j_a)\).
Then their number, with multiplicity, is
\(O_{\tau,\delta,C_0}(B^{1/2+\delta})=o(B)\).
The same estimate holds after points on exact alignments
have been removed.
\end{lemma}

\begin{proof}
Choose a fixed degree \(d\) such that, for
\(M=(d+1)(d+2)/2\), \((M-1)\tau>d\).
Let \(v_d(s,t)\) list the \(M\) monomials of total degree
at most \(d\), including \(1\).
If \(M\) of the vectors \(v_d(i_a,j_a)\) form a basis,
replace each row of that basis in turn by \(v_d(-x,-y)\).
Subtracting the new row from the other rows shows that
the resulting determinant is divisible by the product
of their \(M-1\) pairwise coprime divisors.
Its absolute value is at most
\[
 C_dN^dB^{d(M-1)}<N^{(M-1)\tau}.
\]
All the determinants vanish.
The vector \(v_d(-x,-y)\) would then belong to the
intersection of the \(M\) coordinate hyperplanes of
this basis, which is \(\{0\}\). Its constant coordinate
is one, a contradiction.

Thus a nonzero polynomial \(F(S,T)\) of degree at most
\(d\) vanishes on all the offset points.
For a geometrically integral component of degree at least
two, Pila's uniform estimate
\cite[Theorem~A]{pila1995} gives
\(O_d(B^{1/2+\delta})\) integer points of height \(O(B)\).
A rationally irreducible component that is not geometrically
integral has only \(O_d(1)\) rational points, by intersecting
two conjugate components.

A line through at least two distinct grid points has a
primitive equation \(aS+bT=c\) with
\(|a|,|b|\ll B\), \(|c|\ll B^2\).
Every group on it divides \(ax+by+c\).
This integer is nonzero by the hypothesis, and has size
\(O(NB)\). Coprimality bounds the number of these groups
by \(O_\tau(1)\). A line through only one grid point
needs no separate point count.
Finally, all groups at one point divide its true value
\(x+i_a=O(N)\); their multiplicity is \(O_\tau(1)\).
There are at most \(d\) components. These observations
prove the bound.
\end{proof}

The groups in this lemma can be composite. Only divisibility
and coprimality are used, and the auxiliary curve is not
enumerated. Its degree is chosen after \(\tau\) and before
\(N\).

\paragraph{An anchored small-kernel witness at fixed arity.}
\label{p18:conf:anchored}
Fix integers $r\ge2$, $t\ge2$ and real numbers $\eta,\tau>0$.
Let $M$ be a binary space of physical square relations on $r$ true
windows with centres in $I_N$: its words select an even set in each
window and their total product has even prime valuations.
Write $\ell_v\in M^*$ for the coordinate forms and call a position
active when $\ell_v\ne0$. Designate one window as the anchor.
Assume that:
\begin{itemize}
\item $\dim M\ge\eta L$ and the physical projection into the anchor
is injective;
\item no nonzero form on $M$ has $t$ representations by pairwise
disjoint coordinate sets of size between $1$ and $r-1$;
\item between the anchor and each other block, no two distinct pairs
of active offsets have the same ratio of true values.
\end{itemize}
For every cutoff $H\ge B=L+1$, the anchor then contains two distinct
values $n_1,n_2$ with $K_H(n_i)\le N^{D\tau}$, where
\[
 D=\sum_{j=1}^{r-1}j!(t-1)^{j-1}2^{j(t-1)}.
\]
One threshold, chosen after $r,t,C,\eta,\tau$, precedes the centres,
$M$ and $H$. No injectivity or nonalignment between the other
blocks is assumed.

\begin{proof}
The recovery-exclusion hypothesis gives multiplicity at most $t-1$
for every nonzero coordinate label, by considering singleton recovery
sets. The final counting argument of Lemma~\ref{p1:fixed:recovery}
also bounds the independent $j$-label representations of a nonzero
label by $j!(t-1)^{j-1}$: a maximal disjoint family of their supports
has fewer than $t$ members; its union, of size at most $j(t-1)$,
meets every representation, and induction removes a label. A zero
remaining sum is excluded by independence. Permitting every subset
of the $j$ label fibres gives at most $2^{j(t-1)}$ physical lifts.
Thus at most $D$ physical supports of size at most $r-1$ can
represent a fixed active anchor label.

A prime $\ell>H$ has at most one odd occurrence in each true window.
If it occurs at an active anchor position $v$, its parity row on
$M$ expresses $\ell_v$ as a sum of at most $r-1$ other labels.
Remove zero labels and choose a minimal subrepresentation.
Its labels are distinct, nonzero and independent: otherwise a
nonempty zero-sum subset could be removed without changing the
represented nonzero form. Choose their actual positions in the
prime row. The resulting support $A$ consists of $v$ and positions
in other blocks and uses only active coordinates.

Assign every prime with an active anchor occurrence to one such
support, and let $g_A$ be the product assigned to $A$. Distinct
nontrivial groups are coprime and divide every true value selected
by their support. Every anchor position meets at most $D$ groups,
and its entire odd kernel above $H$ divides their product.
Assign each group to one other block it meets.
For each of the $r-1$ choices, Lemma~\ref{p13:conf:groups} gives
$o(B)$ groups larger than $N^\tau$, with their full multiplicities.
Its no-line hypothesis is exactly the stated active-offset
condition. Equal projected points are covered by the multiplicity
bound in that lemma.

Impose zero only on the anchor positions touched by these large
groups. Each group has a unique anchor position, so the dimension
cost is at most the number of such groups, namely $o(B)$.
A nonzero space remains for large $N$ and its anchor projection
is still injective. A nonzero even projected word uses at least
two positions. Each of those positions was active on the original
space $M$ and met only groups at most $N^\tau$, so its odd kernel
is at most $N^{D\tau}$. The incidence bound is used on $M$ before
erasure; no recovery condition or prime-row classification on the
erased space is assumed.
\end{proof}

\paragraph{Dense primitive hosts and their exact weights.}
\label{p18:conf:fixed-primitive}
For each fixed $r\ge3$, $\eta>0$ and $\varepsilon>0$, primitive
families of $r$ true windows with
$\rho_{\mathrm{rel}}\ge\eta L$ have $O(N^\varepsilon)$ hosts.
Indeed, every effective relation-space projection is injective. Choose
linear sections of the local even-selection maps to lift that space
without increasing its dimension or adding local square kernels.
Lemma~\ref{p1:fixed:recovery}, with $d=r-1$ and $t$ chosen after
$r,\eta$, loses at most $2rB\,t^{-1/(r-1)}\le\eta L/4$ dimensions.
The surviving physical space has dimension at least $3\eta L/4$,
retains all injective projections, and has the required recovery
exclusion. Two exact
units between any pair would give equal nonzero even square classes
in those two blocks: the selected two-value products are nonsquares
by Lemma~\ref{p1:arith:runge}. This contradicts primitivity, since
the pair is a proper subfamily. Thus the active-offset condition
holds for each choice of anchor. Apply the witness with $H=B$ and
density parameter $\eta/2$ at each of the $r$ windows, with the same
fixed constant $D$.
Every centre lies in $\mathcal X_{2D\tau}$.
Lemma~\ref{p13:half:small-kernels} therefore counts at most
$N^{2rD\tau+r\varepsilon'}$ unordered families. Choose $\tau$ and
$\varepsilon'$ after $r,\eta,\varepsilon$ so that this exponent is
at most $\varepsilon$. No auxiliary support or chosen subspace is
enumerated in this last count.

The conversion of this host count to homogeneous probabilities and the
three original weights is given in Proposition~\ref{p13:transfer:primitive}.
That proposition uses the host count proved here, together with the
conditional half-count, and treats the weight comparison at fixed arity.

\begin{theorem}[Dense primitive triples]
\label{p13:conf:dense-primitive}
For fixed \(\eta,\varepsilon>0\), primitive triples of
true windows with \(\rho_{\mathrm{rel}}(T)\ge\eta B\)
have \(O_{\eta,C,\varepsilon}(N^\varepsilon)\) hosts.
Their homogeneous mass is \(O(N^{-2+\varepsilon})\).
The total absolute centered activity of all primitive
triples is \(O(N^{-3/2+\varepsilon})\).
\end{theorem}

\begin{proof}
Apply Proposition~\ref{p13:transfer:primitive} with $r=3$.
The dense condition $\rho_{\mathrm{rel}}(T)\ge\eta B$ implies
$\rho_{\mathrm{rel}}(T)\ge\eta L$, giving the host count and the
homogeneous mass $N^{-2+\varepsilon}$. The global absolute-centered
bound is $N^{-3/2+\varepsilon}$.
\end{proof}

Here homogeneous mass is the sum of the probabilities of the
homogeneous events, with zero right-hand side in place of the affine
start prescription.
It must be distinguished from a weight using the original marginals
$p_x$, which can vanish. For instance, a start at $x=3$ of length
$L=7$ is impossible: the square $4$ forces the common bit on
$3,\ldots,9$ to be zero, and $6=2\cdot3$ then forces the bit at $2$
to be zero, contradicting the left-boundary equation. The corresponding
homogeneous event is compatible and has positive probability.

\begin{corollary}[Dense nonaligned pairs]
\label{p13:conf:dense-pairs}
For fixed \(\eta,\varepsilon>0\), nonaligned pairs
with \(d_{xy}\ge\eta B\) have \(O(N^\varepsilon)\)
hosts.
\end{corollary}

\begin{proof}
Realize their effective relation space with injective
projections. Regularize with
Lemma~\ref{p1:fixed:recovery}, now with \(d=1\),
losing at most \(\eta B/4\).
Every active prime row has two equal labels.
A coordinate belongs to at most \(K=t-1\) active supports.
Confinement bounds the groups above \(N^\tau\) by \(o(B)\).
Kill both endpoints of every such group.
The remaining positive-dimensional space has two active
positions in each block, with kernels at most \(N^{K\tau}\).
The centres lie in \(\mathcal X_{2K\tau}\).
Their count is at most \(N^{4K\tau+2\varepsilon'}\);
choose \(\tau,\varepsilon'\) sufficiently small.
\end{proof}


\paragraph{Completion of the quasi-linear pair profile.}
\label{p18:conf:pair-profile}
We now prove the stronger bounds in Proposition~\ref{p1:arith:hosts},
retaining its full arithmetic nullity
$\rho(x,y)=2L-\dim(S_x+S_y)$.
If both windows are regular, positive $\rho$ is exactly a nonzero
effective intersection; Corollary~\ref{p13:half:pairs} counts these
pairs by $N^{1+\varepsilon'}$. A locally deficient coordinate gives
at most $2N\,|I_N\setminus G_N|=N^{1+o(1)}$ further ordered pairs.
This proves the support bound, even before imposing separation.

For the weighted nonaligned sum, split at
$d_{xy}=\eta L$ with fixed $\eta>0$. Local rigidity gives, uniformly,
\[
 \rho(x,y)=d_{xy}+(L-d_x)+(L-d_y)\le d_{xy}+2,
 \qquad \rho(x,y)\le L+2.
\]
In the low sector, the just proved support bound and
$2^{\eta L+2}\ll_C N^\eta$ give
$N^{1+\eta+\varepsilon'}$. In the high sector,
Corollary~\ref{p13:conf:dense-pairs}, with a smaller fixed density
parameter if necessary to replace $L$ by $B$, gives
$N^{\varepsilon'}$ hosts, each of weight $O_C(N)$.
Choose $\eta,\varepsilon'$ after the final $\varepsilon$.
This proves both assertions with $N^{1+\varepsilon}$ and with a
common threshold before all pairs. Only conditional half-counts,
local rigidity and confinement were used; the pair-profile theorem
and its elementary $6/5$ proof are not inputs to this completion.

\subsubsection{Proof of Proposition~\ref{p13:conf:near-primitive}}
\label{p14:comp:confinement:near}

\begin{proof}
The kernel of each projection of the relation space
has dimension at most \(\theta B\).
Intersect three complements to these kernels.
The resulting subspace loses at most \(3\theta B\)
dimensions and has all projections injective.
Use linear sections of the local even-selection maps to realize
this subspace in physical coordinates, preserving its dimension
and the injectivity of its block projections.

In a given pair, two exact units impose a reduced rational
alignment of height \(O(B)\). Two distinct such alignments
would bound the centres by \(O(B^3)\), which is impossible
at the present scale. The alignment is therefore unique.
If it has \(u\) units, their even-selection space has
dimension \(u-1\); its image loses at most one more
dimension by Lemma~\ref{p1:limits:local-rigidity}.
Thus \(u\le d_{xy}+2\le\theta B+2\).
Kill all the positions of these units in the three pairs,
at a cost of at most \(6\theta B+12\) dimensions.
The remaining space still has macroscopic dimension
and injective projections.

Apply Lemma~\ref{p1:fixed:recovery} with \(d=2\), choosing a fixed
integer \(t\ge2\) such that \(6t^{-1/2}\le\eta/8\).
There are at most \(3B\) physical coordinates, so this regularization
loses at most \(\eta B/8\) further dimensions. Denote the resulting
physical relation space by \(M\). For all sufficiently large \(N\),
\[
 \dim M\ge
 \left(\frac{7\eta}{8}-9\theta\right)B-12
 \ge\frac{\eta L}{2},
\]
since \(\theta<\eta/100\).

Apply the
\hyperref[p18:conf:anchored]{anchored small-kernel witness}
of Section~\ref{p14:comp:confinement} separately to each of the
three windows, using the same space \(M\), the fixed cutoff \(H\),
and density parameter \(\eta/2\).
Its block projections are injective, its coordinate labels satisfy
the recovery exclusion, and its active offsets have no remaining
exact alignments. Let \(D_*=D(3,t)\) be the support-incidence
constant in that witness. Each centre therefore belongs to
\(\mathcal X_{2D_*\tau}\), and
Lemma~\ref{p13:half:small-kernels} bounds the number of triples by
\[
 O\!\left(N^{6D_*\tau+3\varepsilon'}\right).
\]
Choose \(\tau,\varepsilon'>0\) after the requested \(\varepsilon\)
so that \(6D_*\tau+3\varepsilon'\le\varepsilon/2\).
In each application, the incidence bound is used before the final
anchor-coordinate erasure. No auxiliary subspace or deletion
history is enumerated.

Finally,
\(\rho_{\mathrm{rel}}(T)\le d_z+d_{xy}\le L+\theta B\).
The homogeneous weight is at most
\(N^{-2+\theta+o(1)}\). Absorbing the \(o(1)\) in the remaining
exponent margin proves the mass estimate.
\end{proof}

\subsubsection{Proof of Proposition~\ref{p13:conf:strong-profile}}
\label{p14:comp:confinement:strong}

\begin{proof}
Choose the strong pair \(e=\{x,y\}\). Put
\(\delta_z=\dim(S_z\cap(S_x+S_y))
 =\rho_{\mathrm{rel}}(T)-d_{xy}\).
Connectedness gives \(\delta_z>0\), and compatibility gives
\(q(T)=q(e)p_z2^{\delta_z}\).

For \(0<\delta_z<\eta'L\), the conditional count
in Theorem~\ref{p13:half:conditional} gives
\(N^{1/2+\varepsilon'}\) targets of weight at most
\(N^{-1+\eta'+o(1)}\).
Lemma~\ref{p1:arith:dense}, together with
Corollary~\ref{p13:half:pairs} for the independent
part, bounds the total joint mass of strong nonaligned
pairs by \(N^{-1+\varepsilon'}\).
This branch costs \(N^{-3/2+\eta'+O(\varepsilon')}\).

For \(\delta_z\ge\eta'L\), the effective-space form of
Proposition~\ref{p1:fixed:targets}, explained after
\eqref{p13:exact:kappa3}, gives \(N^{\varepsilon'}\)
targets after the fixed pair.
Corollary~\ref{p13:conf:dense-pairs} also gives
\(N^{\varepsilon'}\) source pairs.
Each triple has \(q(T)\ll N^{-3+\sigma+o(1)}\).
Choose \(\eta'\) and the internal margins after
\(\varepsilon\).
\end{proof}
